
\section{Conclusión individual: Jordan}

Durante mi participación en el proyecto me encargué de trabajar en la
verificación de los datos experimentales y en el análisis comparativo
entre TCP y gRPC. Esta parte me permitió comprender mejor cómo se
pueden evaluar dos mecanismos de comunicación a partir de mediciones
reales y no solamente de sus características teóricas.

Al revisar los resultados, observé que gRPC obtuvo una latencia media
menor y un throughput promedio mayor en las pruebas realizadas con
mensajes cortos y largos. Sin embargo, en la prueba de concurrencia
TCP registró un menor tiempo medio de respuesta. La reducción de la latencia
media de gRPC fue del 49,66\,\% en mensajes
cortos y del 32,55\,\% en mensajes largos frente a TCP. En la prueba de tres
solicitudes concurrentes, las medias fueron 6,335300~ms para TCP y
13,068033~ms para gRPC. Esto me permitió
entender que el rendimiento de un mecanismo no se puede valorar
utilizando una sola métrica, ya que su comportamiento puede variar
según las condiciones de ejecución.

También trabajé en la validación de los archivos CSV y en las pruebas
automatizadas para detectar datos incompletos, valores incorrectos
y resultados estadísticos inconsistentes. Comprobar que las pruebas
se ejecutaran correctamente fue importante para tener mayor confianza
en la información utilizada durante el análisis.

Una de las principales limitaciones que identifiqué fue que los tres
nodos se ejecutaron en el mismo equipo, utilizando la interfaz local.
Además, la prueba de concurrencia se limitó a tres solicitudes, por
lo que todavía sería necesario realizar experimentos adicionales
en diferentes equipos y con una mayor cantidad de comunicaciones.

En general, este trabajo me ayudó a relacionar los conceptos de
comunicación distribuida con su aplicación práctica. También aprendí
la importancia de validar los datos antes de interpretarlos y de
presentar conclusiones que realmente correspondan a los resultados
obtenidos.
